% Section 7 — Limitations and future research
\section{Limitations and Future Research}
\label{sec:limits}

\subsection{Dataset scope}
\label{sec:l-data}
Two retail environments, however contrasted, are still two. Both are
brick-and-mortar-style daily demand; neither covers e-commerce spikes,
perishables with expiry, promotions-driven regimes, or supply-side
censoring beyond the lost-sales assumption.

\subsection{Series selection}
\label{sec:l-series}
M5 results describe the stratified 500, not the 30,490-series catalog.
Stratification protects archetype coverage but cannot guarantee catalog-level
point estimates; scaling the full ladder to the catalog is compute-bound
(SARIMA alone cost 2,382\,s for 500 series).

\subsection{Forecast horizon}
\label{sec:l-horizon}
All conclusions are conditional on $h=28$. Shorter horizons favor naive and
averaging methods; longer ones may favor trend/seasonal machinery. The
robustness grid varies policy, not horizon.

\subsection{Inventory-policy scope}
\label{sec:l-policy}
One policy family (daily-review order-up-to, lost sales, no capacity, no
minimum order quantities, no perishability, stationary costs) carries all
decision conclusions. Backordering regimes, batch constraints, or stochastic
lead times could reorder the boards --- which is precisely why the
sensitivity discipline, not any single ranking, is the transferable asset.

\subsection{Feature scope}
\label{sec:l-features}
Univariate demand only: no prices, promotions, SNAP/calendar events, or
cross-series covariates. The LSTM's advantage could widen with features; the
classical methods' standing could improve with regressors. Both are open
questions, not assumptions.

\subsection{Archetype limitations}
\label{sec:smooth-cap}
The Smooth archetype contains exactly one series ($n=1$). Nothing in this
report generalizes from it, and any Smooth figure is shown for completeness
only.

\subsection{SARIMA scope}
\label{sec:l-sarima}
SARIMA is Store-only by design; no claim is made about seasonal
autoregression on intermittent demand. The 100-series subset is exploratory
history, never the headline.

\subsection{Computational considerations}
\label{sec:l-compute}
Per-fit costs ($\approx$0.03\,s ARIMA, $\approx$0.6\,s SARIMA, plus LSTM
training) bound full-catalog scaling. Per-series ARIMA order search was
rejected on this ground; future work with larger budgets could revisit it.

\subsection{LLM experiment --- future work only}
\label{sec:llm-future}
The research program's final rung --- local LLM forecasting (controlled vs.\
context-enhanced prompting, Ollama backend, leak-guarded prompt builder in
\texttt{11\_src/llm\_experiment\_design.md}) --- is designed but deliberately
unexecuted, held behind a frozen benchmark so the benchmark could not be tuned
against it. \textbf{No LLM result exists; none is presented, implied, or
charted anywhere in this report.} Executing that phase against the frozen
evidence base is the natural next step.

\subsection{Proposed future extensions}
\label{sec:l-future}
In priority order: (1) execute the LLM phase; (2) add price/promotion/calendar
covariates; (3) widen the policy family (backorders, MOQs, perishability);
(4) vary the horizon (7/56 days); (5) scale M5 toward the full catalog for the
cheap rungs; (6) test additional retail environments beyond the two studied.
